\newpage
\begin{center}
	\textbf{\large 3. Результаты}
\end{center}
\refstepcounter{chapter}
\addcontentsline{toc}{chapter}{3. Результаты}

Тестирование разработанного алгоритма локализации выполнялось на наборе изображений с ручной разметкой (Ground Truth), подготовленной в инструменте LabelMe.
Для каждого изображения алгоритм либо возвращал четырёхугольник области QR-кода (четыре вершины в порядке \texttt{TL--TR--BR--BL}), либо сообщал об отсутствии результата.
Далее результат сравнивался с Ground Truth по метрикам IoU и ошибкам по углам.

Для оценки точности локализации использовались следующие метрики:
\begin{itemize}
	\item \textbf{IoU (Intersection over Union)} --- отношение площади пересечения области, найденной алгоритмом, к площади объединения найденной области и эталонной разметки.
	      Значение $IoU \in [0,1]$, где 1 соответствует полному совпадению.
	\item \textbf{Средняя ошибка по углам} (\texttt{corner\_err\_mean\_px}) --- среднее евклидово расстояние (в пикселях) между соответствующими вершинами предсказанного многоугольника и разметки:
	      \[
		      e_{mean}=\frac{1}{4}\sum_{i=1}^{4} \|q_i - q_i^{GT}\|.
	      \]
	\item \textbf{Максимальная ошибка по углам} (\texttt{corner\_err\_max\_px}) --- максимальная ошибка среди четырёх вершин:
	      \[
		      e_{max}=\max_{i\in\{1,2,3,4\}} \|q_i - q_i^{GT}\|.
	      \]
\end{itemize}
Метрики по углам позволяют дополнительно оценить характер ошибки (например, наличие одного сильно ``съехавшего'' угла), даже если IoU остаётся относительно высоким.
Результаты тестирования доступны в файле \texttt{results.csv}.

\subsubsection{Сводные результаты}
По приведённой выборке (28 изображений, на всех алгоритм вернул результат, \texttt{algo\_found = YES}) получены следующие значения:
\begin{itemize}
	\item \textbf{Средний IoU:} $\approx 0.935$.
	\item \textbf{Медианный IoU:} $\approx 0.964$.
	\item \textbf{Минимальный IoU:} $0.784843$ (файл \texttt{test22.jpg}).
	\item \textbf{Максимальный IoU:} $0.988239$ (файл \texttt{test28.jpg}).
\end{itemize}

\begin{figure}[!h]
	\centering
	\includegraphics[width=0.45\textwidth]{im06.png}
	\caption{Файл \texttt{test22.jpg}: увеличивается влияние перспективы}
	\label{fig:a1}
\end{figure}

\begin{figure}[!h]
	\begin{center}
		\includegraphics[width=250pt]{im07.png}
		\caption{Файл \texttt{test28.jpg}: QR-код практически перпендикулярен камере}
		\label{fig:a2}
	\end{center}
\end{figure}

\begin{figure}[!h]
	\begin{center}
		\includegraphics[width=250pt]{im08.png}
		\caption{Файл \texttt{test19.jpg}: повёрнутые QR-коды успешно детектируются}
		\label{fig:a3}
	\end{center}
\end{figure}

\newpage

По ошибкам углов:
\begin{itemize}
	\item \textbf{Средняя ошибка по углам:} $\approx 10.8$ px.
	\item \textbf{Медианная ошибка по углам:} $\approx 7.0$ px.
	\item \textbf{Средняя максимальная ошибка:} $\approx 17.1$ px.
	\item \textbf{Медианная максимальная ошибка:} $\approx 12.6$ px.
\end{itemize}

Полученные значения подтверждают, что в большинстве случаев границы QR-кода определяются с высокой точностью: для значительной части изображений IoU превышает $0.95$, а средняя ошибка по углам составляет единицы пикселей даже при разрешении изображения порядка нескольких тысяч пикселей.

\subsubsection{Анализ по группам качества}
Для более содержательного анализа результаты были разделены на группы по значению IoU.

\paragraph{Высокое качество (IoU $\ge 0.95$).}
К этой группе относится большинство примеров (например, \texttt{test0.jpg}, \texttt{test11--test17.jpg}, \texttt{test25.jpg}, \texttt{test35.jpg}).
Для таких изображений средняя ошибка углов, как правило, не превышает $5$--$8$ пикселей, а максимальная ошибка --- порядка $10$--$15$ пикселей.
Это соответствует точной оценке масштаба и положения QR-кода, а также корректному выбору тройки finder patterns.

\begin{figure}[!h]
	\centering
	\includegraphics[width=0.5\textwidth]{im10.png}
	\caption{Файл \texttt{test0.jpg}: высокое качество}
	\label{fig:a4}
\end{figure}

\paragraph{Среднее качество ($0.90 \le$ IoU $< 0.95$).}
В эту группу входят примеры \texttt{test6.jpg}, \texttt{test9.jpg}, \texttt{test10.jpg}, \texttt{test20.jpg}, \texttt{test29.jpg}.
Характерно увеличение ошибок по углам (например, средняя ошибка 11--15 px и максимальная до 25--30 px).
На практике это может соответствовать небольшому перспективному искажению, локальному повороту, либо неточности оценки размера finder pattern (что влияет на восстановление внешней границы по геометрической модели).

\begin{figure}[!h]
	\centering
	\includegraphics[width=0.5\textwidth]{im11.png}
	\caption{Файл \texttt{test6.jpg}: среднее качество}
	\label{fig:a5}
\end{figure}

\paragraph{Низкое качество (IoU $< 0.90$).}
К этой группе относятся:
\texttt{test22.jpg} (IoU $0.784843$),
\texttt{test26.jpg} (IoU $0.872138$),
\texttt{test30.jpg} (IoU $0.851586$),
\texttt{test31.jpg} (IoU $0.847256$),
\texttt{test32.jpg} (IoU $0.804507$).
Для этих случаев характерны большие ошибки углов (средняя ошибка до $38$ px и максимальная до $52$ px).
Практически это означает, что хотя алгоритм нашёл QR-код, один или несколько углов определены заметно неточно.
Основные причины:
\begin{itemize}
	\item ошибочная оценка размера finder pattern (например, был измерен не внешний, а внутренний квадрат маркера);
	\item влияние перспективы/наклона, при котором аффинная аппроксимация границ становится недостаточной;
	\item деградация бинаризации на отдельных кадрах (частичное размытие, локальные блики), что ухудшает качество контуров и геометрических оценок.
\end{itemize}


\begin{figure}[!h]
	\centering
	\includegraphics[width=0.5\textwidth]{im12.png}
	\caption{Файл \texttt{test30.jpg}: низкое качество}
	\label{fig:a6}
\end{figure}

\begin{figure}[!h]
	\centering
	\includegraphics[width=0.5\textwidth]{im13.png}
	\caption{Файл \texttt{test32.jpg}: низкое качество}
	\label{fig:a6}
\end{figure}

В результате выполнения работы было разработано консольное приложение для локализации QR‑кода на изображениях большого разрешения.
Алгоритм реализует полный цикл обработки: предобработку изображения (перевод в оттенки серого, сглаживание, бинаризация), поиск и фильтрацию кандидатов finder patterns по контурам и их иерархии, выбор наиболее вероятной тройки маркеров по геометрическому score и восстановление границ QR‑кода в виде четырёхугольника.
Для проверки качества была подготовлена эталонная разметка (Ground Truth) и реализовано сравнение результатов с ней по количественным метрикам (IoU и ошибки по углам), что позволило выявлять характерные ошибки и уточнять параметры.
